\subsection{COMPACT LIE ALGEBRAS}

PROPOSITION 1. Let \(\mathfrak{g}\) be a (real) Lie algebra. The following conditions are equivalent:

\begin{enumerate}
\def\labelenumi{(\roman{enumi})}
\item
  g is isomorphic to the Lie algebra of a compact Lie group.
\item
  The group \(\operatorname{Int}(\mathfrak{g})\) (Chap. III, §6, no. 2, Def. 2) is compact.
\item
  \(\mathfrak{g}\) has an invariant bilinear form (Chap. I, §3, no. 6) that is symmetric, positive and separating.
\item
  g is reductive (Chap. I, §6, no. 4, Def. 4); for all \(x \in g\) , the endomorphism ad x is semi-simple, with purely imaginary eigenvalues.
\item
  \(\mathfrak{g}\) is reductive and its Killing form B is negative.
\item
  \(\Longrightarrow\) (ii): if g is the Lie algebra of a compact Lie group G, the group \(\operatorname{Int}(\mathfrak{g})\) is separating and isomorphic to a quotient of the compact group \(G_{0}\) (Chap. III, §6, no. 4, Cor. 4), hence is compact.
\item
  \(\Longrightarrow\) (iii): if the group \(\operatorname{Int}(\mathfrak{g})\) is compact, there exists a symmetric bilinear form on g that is positive, separating and invariant under \(\operatorname{Int}(\mathfrak{g})\) (no. 1), hence also invariant under the adjoint representation of g.
\item
  \(\Longrightarrow\) (iv): if (iii) is satisfied, the adjoint representation of g is semi-simple (no. 1), hence g is reductive; moreover, the endomorphisms ad x, for \(x \in g\) , have the indicated properties (no. 1).
\item
  \(\Longrightarrow\) (v): for all \(x\in \mathfrak{g}\) , \(\mathrm{B}(x,x) = \mathrm{Tr}((\mathrm{ad}x)^2)\) ; consequently, \(\mathrm{B}(x,x)\) is the sum of the squares of the eigenvalues of ad \(x\) , and hence is negative if these are purely imaginary.
\item
  \(\Longrightarrow\) (i): assume that g is reductive, hence the product of a commutative subalgebra c and a semi-simple subalgebra s (Chap. I, §6, no. 4, Prop. 5). The Killing form of s is the restriction of the form B to s, hence is negative and separating if B is negative. The subgroup \(\operatorname{Int}(\mathfrak{s})\) of \(\mathbf{GL}(\mathfrak{s})\) is closed (it is the identity component of \(\operatorname{Aut}(\mathfrak{s})\) , Chap. III, §10, no. 2, Cor. 2) and leaves the separating positive form -B invariant; thus, it is compact, and s is isomorphic to the Lie algebra of the compact Lie group \(\operatorname{Int}(\mathfrak{s})\) . Further, since c is commutative, it is isomorphic to the Lie algebra of a torus T. Thus g is isomorphic to the Lie algebra of the compact Lie group \(\operatorname{Int}(\mathfrak{s}) \times \operatorname{T}\) .
\end{enumerate}

DEFINITION 1. A compact Lie algebra \(^{4}\) is a Lie algebra that has properties (i) to (v) of Proposition 1.

Thus, the compact Lie algebras are the products of a commutative algebra with a compact semi-simple algebra. In other words, a Lie algebra is compact if and only if it is reductive and its derived Lie algebra is compact.

The Lie algebra of a compact Lie group is compact.

PROPOSITION 2. a) The product of a finite number of Lie algebras is a compact Lie algebra if and only if each factor is compact.

\begin{enumerate}
\def\labelenumi{\alph{enumi})}
\setcounter{enumi}{1}
\item
  A subalgebra of a compact Lie algebra is compact.
\item
  Let \(\mathfrak{h}\) be an ideal of a compact Lie algebra \(\mathfrak{g}\) . Then the algebra \(\mathfrak{g} / \mathfrak{h}\) is compact and the extension \(\mathfrak{h} \to \mathfrak{g} \to \mathfrak{g} / \mathfrak{h}\) is trivial.
\end{enumerate}

Assertions \(a\) and \(b\) follow from the characterization (iii) of Prop. 1. Part \(c\) follows from \(a\) and the fact that, in a reductive Lie algebra, every ideal is a direct factor (Chap. I, §6, no. 4, Cor. of Prop. 5).

PROPOSITION 3. Let G be a Lie group of which the group of connected components is finite. The following conditions are equivalent:

\begin{enumerate}
\def\labelenumi{(\roman{enumi})}
\item
  The Lie algebra \(\mathrm{L}(\mathrm{G})\) is compact.
\item
  The group \(\mathrm{Ad}(\mathrm{G})\) is compact.
\item
  There exists a separating positive symmetric bilinear form on L(G) invariant under the adjoint representation of G.
\end{enumerate}

*(iv) G has a riemannian metric invariant under left and right translations.*

\begin{enumerate}
\def\labelenumi{(\roman{enumi})}
\item
  \(\Longrightarrow\) (ii): if L(G) is compact, the group \(\mathrm{Ad}(\mathrm{G}_{0}) = \mathrm{Int}(\mathrm{L}(\mathrm{G}))\) is compact; since it has finite index in \(\mathrm{Ad}(\mathrm{G})\) , this latter group is also compact.
\item
  \(\Longrightarrow\) (iii): this follows from no. 1.
\item
  \(\Longrightarrow\) (i): since \(\operatorname{Int}(\mathrm{L}(\mathrm{G})) \subset \operatorname{Ad}(\mathrm{G})\) , this follows from the characterization (iii) of Prop. 1.
\end{enumerate}

\(^{*}\) (iii) \(\Longleftrightarrow\) (iv): this follows from Chap. III, §3, no. 13.*

